\documentclass[10pt,a4paper]{amsart}
\usepackage[margin=19mm]{geometry}
\usepackage{amsmath,amssymb,mathtools}
\usepackage[scr=rsfs,cal=euler]{mathalfa}
\usepackage{xfrac}
\usepackage[unicode,hidelinks,bookmarks=false]{hyperref}
\newcommand{\ie}{i.e.\ }
\newcommand{\cf}{cf.\ }
\newcommand{\resp}{respectively\ }
\newcommand{\Subsection}{\subsection}
\providecommand{\nobreakdash}{\nobreak\hbox{-}}
\setlength{\emergencystretch}{2em}
\multlinegap0pt
\usepackage{bm}
\usepackage{bbm}
\usepackage{dsfont}
\usepackage{xspace}
\usepackage{cancel}
\usepackage{mathtools}
\usepackage{amsthm}
\makeatletter
\@ifundefined{theorem}{\newtheorem{theorem}{Theorem}}{}
\@ifundefined{proposition}{\newtheorem{proposition}[theorem]{Proposition}}{}
\makeatother
\newcommand{\Hyp}{\mathrm{\textbf{Hyp}}}
\newcommand{\MHyp}{\mathrm{\textbf{MHyp}}}
\newcommand{\Simp}{\mathrm{\textbf{Simp}}}
\renewcommand{\H}{\mathcal{H}}
\title[A survey of simplicial, relative, and chain complex homology]{A survey of simplicial, relative, and chain complex homology theories for hypergraphs}
\author{Ellen Gasparovic, Emilie Purvine, Radmila Sazdanovic, Bei Wang, Yusu Wang, Lori Ziegelmeier}
\date{Source: arXiv:2409.18310v2 (CC BY 4.0)}
\begin{document}
\maketitle

\noindent\emph{Source and licence.} A survey of simplicial, relative, and chain complex homology theories for hypergraphs. Version arXiv:2409.18310v2 (CC BY 4.0), distributed under the Creative Commons Attribution 4.0 International licence, as recorded in the arXiv licence metadata for this exact version and shown on the abstract page https://arxiv.org/abs/2409.18310v2. Full rights notice: licenses/arxiv-2409.18310-CC-BY-4.0.txt.\par\smallskip

\noindent\textit{Editorial scope.} Counted unit: the Proposition whose source label is \texttt{prop:closure\_functor} in the article (source0.tex lines 622--624, source label \texttt{prop:closure\_functor}), delivered together with the article's own complete original proof of it (source0.tex lines 625--649, one \texttt{proof} environment of 25 lines). No mathematics is written, repaired or re-derived: statement and proof are the article's own text line for line. Required context retained uncounted: none. Every definition, display and label of those lines is retained unchanged. Omitted from this package and not redistributed: 1--621, 650--2322. Nothing inside a retained excerpt is omitted or altered: every retained body is delivered exactly as the article's lines are written, so no omission receipt of any kind accompanies this package. Citations: the retained excerpts carry no citation command at all, so no bibliography entry of the article is transcribed and no reference is redistributed. Reference adaptation: none was needed and none was made; every reference inside the delivered unit resolves inside it, so no unresolved reference or citation is produced. Printed slips kept literal: no printed word, symbol, hypothesis or number of the counted statement or of its proof was altered. Layout and class: the article's own class is \texttt{amsart} with packages including inputenc, xcolor, amsfonts, amssymb, graphicx, tikz, tikz-cd, fancyhdr, bbm, subcaption, footmisc, hyperref; the wrapper is the lane's pinned \texttt{amsart} editorial wrapper at 10pt on A4, which restates the theorem-like environments the retained text needs and adds the article's own macro definitions that the retained text needs (\texttt{\textbackslash H}, \texttt{\textbackslash Hyp}, \texttt{\textbackslash MHyp}, \texttt{\textbackslash Simp}), with extra wrapper packages (bm, bbm, dsfont, xspace, cancel, mathtools). The delivered numbering is the unit's own. Assets: no retained span contains a figure environment, a graphics-inclusion command, a PDF-page inclusion, a table or a TikZ picture; no image file is copied into this package, the package declares no asset dependency and no image file is redistributed with it. Licence: version arXiv:2409.18310v2 of this article is distributed under the Creative Commons Attribution 4.0 International licence, as recorded in the arXiv licence metadata for this exact version and shown on the abstract page https://arxiv.org/abs/2409.18310v2. A full rights notice is delivered at licenses/arxiv-2409.18310-CC-BY-4.0.txt. Characters: every retained excerpt is pure ASCII, so the delivered engine, XeLaTeX with its default Latin Modern text font, reports no missing character.
\begin{proposition}\label{prop:closure_functor}
The hypergraph upper closure operation is a functor $\Delta : \MHyp \rightarrow \Simp$.
\end{proposition}

\begin{proof}
The first observation is that $\Delta$ is a composition of the collapse functor, $Col:\MHyp\rightarrow\Hyp$, and the process of adding in sub-edges, which we have not yet formally defined as an object map and morphism map. We do so now for $\Hyp$ to $\Simp$ and prove that it is a functor. Then its composition with $Col$ will also be a functor.
\begin{samepage}
\begin{align*}
    \Delta : \Hyp& \rightarrow \Simp\\
        (V, E)& \mapsto (V, \{\emptyset \neq \sigma \subseteq V : \exists e \in E \text{ with }\sigma \subseteq e\}) \hspace{2em} \textrm{(object map)}\\
        f& \mapsto f \hspace{17em} \textrm{(morphism map)}
\end{align*}
\end{samepage}

The fact that $\Delta(Id_\Hyp) = Id_\Simp$ and $\Delta(f_1\circ f_2) = \Delta(f_1)\circ\Delta(f_2)$ are true by definition since the morphism map is the identity. To show $\Delta$ is a functor it is left to show that if $f$ is a morphism in $\Hyp$ from $\H_1=(V_1, E_1)$ to $\H_2=(V_2, E_2)$ then $\Delta(f)$ is a simplicial map from $\Delta(\H_1)$ to $\Delta(\H_2)$. 

Let $\sigma \in \Delta(\H_1)$, then there is an $e \in E_1$ such that $\sigma \subseteq e$. Since $f$ is a morphism in $\Hyp$ it must be that $\{f(v)\}_{v \in e}$ is an edge in $\H_2$. And of course $\{f(v)\}_{v \in \sigma}$ is a subset of this edge.
Note that since $\{f(v)\}_{v \in e}$ is an edge in $\H_2$ it is also a simplex in $\Delta(\H_2)$. Finally since any subset of a simplex is also a simplex, namely $\{f(v)\}_{v \in \sigma}$, $f$ is a simplicial map and we have the desired result. 


    %\cEP{old}
    %The object map, $\Delta(\H)$, was defined earlier but we have not yet defined the map on morphisms. Let $(f, g)$ be a morphism in $\Hyp$ from $\H_1=(V_1, E_1, \epsilon_1)$ to $\H_2 = (V_2, E_2, \epsilon_2)$. Recall from Section \ref{sec:categories} that $f$ is a map on edges and $g$ is a map on vertices. Then we define $\Delta((f, g)) := g$. The fact that $\Delta(Id_\Hyp) = Id_\Simp$ and $\Delta((f_1, g_1)\circ (f_2, g_2)) = \Delta((f_1, g_1))\circ\Delta((f_2, g_2))$ are inherited from $\Hyp$. To show $\Delta$ is a functor it is left to show that $g$ is a simplicial map from $\Delta(\H_1)$ to $\Delta(\H_2)$. 
    %
    %Let $\sigma \in \Delta(\H_1)$, then there is an $e \in E_1$ such that $\sigma \subseteq \epsilon_1(e)$, and $f(e) \in E_2$. By the definition of morphisms in $\Hyp$ it must be that $\po(g)(\epsilon_1(e)) = \epsilon_2(f(e))$. The left side can be rewritten as 
    %$\po(g)(\{v : v \in \epsilon_1(e)\}) = \{g(v) : v \in \epsilon_1(e)\}$.
    %Since $\sigma \subseteq \epsilon_1(e)$ we can put this all together and see that
    %\[ \{g(v) : v \in \sigma\} \subseteq \epsilon_2(f(e)). \]
    %Note that since $f(e)$ is an edge in $\H_2$, $\epsilon_2(f(e))$ is a simplex in $\Delta(\H_2)$. Finally since any subset of a simplex is also a simplex, $g$ is a simplicial map and we have the desired result.
\end{proof}

\end{document}
